\chapter{Accelerators for Quantitative Work}\label{ch:pl:accelerators-for-quantitative-work}

A bank moved its overnight counterparty-risk Monte Carlo to graphics processors and it ran many times faster. The same bank tried the
same devices for the intraday pricing of single trades, and nothing moved: each request spent longer crossing to the device and back
than computing on it. Both results were predictable from three numbers per workload --- how much arithmetic it does, how many bytes it
moves, and how many bytes must cross to the device --- and two numbers per machine. This chapter measures those numbers on the laptop's
processor, takes a data-centre device's from its product page, and predicts both outcomes before anyone buys anything.

\section{What a graphics processor is good at}

A graphics processor (an accelerator in the sense of Book~12, chapter~23) has thousands of simple arithmetic units running the same
instruction on different data, and memory attached to it that is several times faster than a server's. It is good at the same thing
SIMD instructions are good at on a CPU (Book~13, chapter~14), at a much larger scale: a long loop of identical, independent arithmetic on
arrays already in its memory. It is bad at three things quantitative code often does.

\begin{definition}[Host--device transfer]\label{def:pl:accelerators-for-quantitative-work:transfer}
A \emph{host--device transfer}\index{host--device transfer} moves data between the CPU's memory (the host) and an accelerator's memory
(the device) over a link such as PCI Express; every input a device computation reads and every result it returns crosses it, at a
bandwidth far below either memory's and with a fixed cost per crossing.
\end{definition}

\begin{definition}[Thread divergence]\label{def:pl:accelerators-for-quantitative-work:divergence}
\emph{Thread divergence}\index{thread divergence} happens when threads that execute in lockstep take different branches: the group runs
each branch in turn with the threads not on it idle, so code with data-dependent branches (early exercise, barrier checks, path-dependent
payoffs with many cases) runs at a fraction of the device's rate.
\end{definition}

\begin{omfigure}
\begin{tikzpicture}[font=\scriptsize,box/.style={draw,rounded corners=2pt,align=center,minimum height=0.8cm,inner sep=3pt},>=Stealth]
\node[box,fill=omMeth!15,minimum width=2.6cm] (cpu) at (0,0) {CPU cores};
\node[box,fill=gray!10,minimum width=2.6cm] (hm) at (0,-1.6) {host memory};
\node[box,fill=omDef!15,minimum width=3.2cm] (gpu) at (8.5,0) {device: thousands\\ of arithmetic units};
\node[box,fill=gray!10,minimum width=3.2cm] (dm) at (8.5,-1.6) {device memory};
\draw[<->,thick] (cpu) -- node[left]{24.0 GB/s, one core} (hm);
\draw[<->,very thick] (gpu) -- node[right]{3.35 TB/s} (dm);
\draw[<->,thick,omThm] (hm.east) -- node[above]{PCIe Gen5: 128 GB/s} node[below]{every input and result crosses} (dm.west);
\draw[->,dashed] (cpu.east) -- node[above=2pt]{launch and synchronise: $L$ per request} (gpu.west);
\end{tikzpicture}

\omcaption{Where the time goes when work moves to a device. The device's memory is 140 times faster than what one core of the laptop
reaches, but inputs and results cross a link slower than either memory, and every request pays a fixed launch and synchronisation cost $L$
(the chapter assumes 10~microseconds).}\label{fig:pl:accelerators-for-quantitative-work:path}
\end{omfigure}

The third is size. A computation that takes a microsecond on a CPU cannot be faster on a device whose launch and synchronisation alone
take several: small work, like a single trade's price, pays a fixed cost that large work amortises.

\section{The roofline model}

\begin{definition}[Arithmetic intensity]\label{def:pl:accelerators-for-quantitative-work:intensity}
The \emph{arithmetic intensity}\index{arithmetic intensity} of a kernel is the number of floating-point operations it performs per byte
it moves between memory and the processor.
\end{definition}

\begin{definition}[Roofline model]\label{def:pl:accelerators-for-quantitative-work:roofline}
The \emph{roofline model}\index{roofline model} (Williams, Waterman and Patterson, 2009) bounds the arithmetic rate a kernel of intensity
$I$ can attain on a machine of peak rate $P$ and memory bandwidth $B$ by $\min(P,\ I B)$. Plotted on logarithmic axes it is a slanted
line (bandwidth) meeting a flat one (peak) at the ridge point $I^* = P/B$.
\end{definition}

\begin{definition}[Memory-bound kernel, compute-bound kernel]\label{def:pl:accelerators-for-quantitative-work:bound}
A kernel whose intensity is below the machine's ridge point is \emph{memory-bound}\index{memory-bound kernel}: its rate is limited by
bandwidth, and more arithmetic units do not help it. A kernel above the ridge is \emph{compute-bound}\index{compute-bound kernel}.
\end{definition}

The roofline turns a question about hardware into two numbers about code. The chapter's model (\texttt{firm.roofline}) adds the transfer
and the fixed cost: on a device, a kernel takes at least $L + X/\ell + \max(F/P, M/B)$ seconds, with $F$ its operations, $M$ its memory
bytes, $X$ the bytes that cross a link of bandwidth $\ell$, and $L$ the launch and synchronisation latency of a request.

\omcode{code/firm/roofline/firm_roofline.py}{55}{75}{A roof and the end-to-end time of a kernel on it: the roofline bound, plus a
launch latency and a transfer over the host link for a device.}\label{lst:pl:accelerators-for-quantitative-work:roof}

\section{Measuring the laptop's roof}

One core of this laptop (an Intel Core Ultra 7 155H, under WSL2) was measured: the peak by a $2000\times2000$ double-precision matrix
product on one BLAS thread, the bandwidth by a STREAM-style triad $a_i = b_i + s c_i$ in C++20 on arrays of 128~MB, counting 24 bytes per
element. The core reaches 66.9~GFLOP/s and 24.0~GB/s: its ridge is at 2.79 operations per byte.

\omcode{code/firm/roofline/cpp/firm_roofline_triad.cpp}{11}{22}{The bandwidth measurement: the triad, timed best of ten.}\label{lst:pl:accelerators-for-quantitative-work:triad}

Four reference kernels were timed on the same core (\cref{tab:pl:accelerators-for-quantitative-work:kernels}): a step of geometric
Brownian motion over $2^{23}$ paths ($S \mathrel{*}= e^{a + bZ}$), a basket payoff over 500\,000 paths of 20 assets, the aggregation of a
cube of trade values (2\,000 paths, 50 dates, 200 trades) into an expected-positive-exposure profile, and the matrix product. The operations
are counted by the convention of \texttt{firm.roofline}: additions and multiplications count one, and so does a call to \texttt{exp},
although it costs many cycles.

\begin{table}[ht]
\centering\small
\begin{tabular}{lrrrr}
\toprule
kernel & intensity (flop/byte) & attained (GFLOP/s) & roof (GFLOP/s) & share of roof \\
\midrule
path step & 0.17 & 0.65 & 3.99 & 16\% \\
basket payoff & 0.25 & 4.82 & 5.99 & 81\% \\
exposure aggregation & 0.13 & 2.63 & 3.03 & 87\% \\
matrix product & 166.7 & 64.2 & 66.9 & 96\% \\
\bottomrule
\end{tabular}
\caption{Four kernels on one core of the laptop: intensity by the chapter's counting convention, the rate attained (best of five runs),
the roof at that intensity and the share reached. Data: \texttt{bench\_roofline.py}.}
\label{tab:pl:accelerators-for-quantitative-work:kernels}
\end{table}

Three of the four are \omterm{def:pl:accelerators-for-quantitative-work:bound}{memory-bound}, and the two that are simple loops reach 81 and 87\% of the bandwidth roof: they are done, and only more
bandwidth would speed them up. The path step sits far below its roof because the roofline counts operations and \texttt{exp} is one
operation that takes many cycles; its true limit is the rate of the exponential. The matrix product, the only \omterm{def:pl:accelerators-for-quantitative-work:bound}{compute-bound kernel}, defines
the peak: it reaches 96\% of it here only because the peak is its best over more runs than the table's five.

\begin{omfigure}
\begin{tikzpicture}
\begin{loglogaxis}[width=12cm,height=6.4cm,xlabel={arithmetic intensity (flop per byte, log scale)},ylabel={GFLOP/s (log scale)},xmin=0.01,xmax=1000,
  ymin=0.1,ymax=100000,grid=major,grid style={gray!20},tick label style={font=\small},label style={font=\small},table/col sep=comma]
\addplot[omThm,very thick] table[x=intensity,y=cpu_gflops]{figdata/platforms/14-accelerators-for-quantitative-work/roofline.csv};
\addplot[omDef,very thick,dashed] table[x=intensity,y=device_gflops]{figdata/platforms/14-accelerators-for-quantitative-work/roofline.csv};
\addplot[only marks,mark=*,mark size=2pt,black]
  table[x=intensity,y=gflops]{figdata/platforms/14-accelerators-for-quantitative-work/measured_kernels.csv};
\node[font=\scriptsize,anchor=west] at (axis cs:0.18,0.65) {path step};
\node[font=\scriptsize,anchor=west] at (axis cs:0.29,4.8) {basket payoff};
\node[font=\scriptsize,anchor=south east,align=right,inner sep=1pt] at (axis cs:0.12,2.9) {exposure\\aggregation};
\node[font=\scriptsize,anchor=south east] at (axis cs:167,64) {matrix product};
\node[omThm,font=\scriptsize,anchor=north] at (axis cs:30,50) {one laptop core};
\node[omDef,font=\scriptsize,anchor=north] at (axis cs:100,30000) {data-centre device (FP64)};
\end{loglogaxis}
\end{tikzpicture}

\omcaption{Two rooflines: one core of the laptop (measured) and a data-centre graphics processor in double precision (from its product
page), with the four kernels as measured on the core. The core's ridge is at 2.8 operations per byte, the device's at 10.1: the kernels that
are \omterm{def:pl:accelerators-for-quantitative-work:bound}{memory-bound} on the core are more so on the device. Data: \texttt{bench\_roofline.py}, \texttt{fig\_roofline.py}.}
\label{fig:pl:accelerators-for-quantitative-work:roofline}
\end{omfigure}

\begin{dated}{2026-09}{A data-centre device}\label{dat:pl:accelerators-for-quantitative-work:h100}
NVIDIA's product page gives, for the H100 in its SXM form: 34 teraFLOPS in double precision (67 on its tensor cores), 67 in single, 80~GB of
memory at 3.35~TB/s, a 900~GB/s NVLink interconnect, 128~GB/s over PCIe Gen5, and up to 700~W. Against one core of the laptop it offers 508
times the double-precision peak and 140 times the bandwidth. The chapter uses these figures as the device's roof and its PCIe figure as the
host link, the most favourable reading of it.
\end{dated}

\section{Monte Carlo and risk on a device}

The overnight workload is an exposure Monte Carlo for one netting set: 10\,000 paths of 50 risk factors on 120 monthly dates, 5\,000 trades
valued on every path and date (ten operations each: a linear trade on one factor, discounted), and the values aggregated into an
expected-exposure profile (Book~6, chapters 17 and~18). It performs 66~billion operations and moves 98~GB through memory: \omterm{def:pl:accelerators-for-quantitative-work:bound}{memory-bound} on both
machines. Only the trades' descriptions go to the device (100 bytes each) and only the profile comes back; the paths are generated where
they are used, with a counter-based generator (Book~4, chapter~26) so that the device's paths are the CPU's.

\omcode{code/platforms/14-accelerators-for-quantitative-work/python/pl_roofline.py}{43}{65}{The two workloads in the model: the exposure
job's three kernels, its end-to-end time on the device, and a request to price $n$ trades.}\label{lst:pl:accelerators-for-quantitative-work:model}

On the core, at its roof, the job takes 4.09~seconds; on the device, 29.3~milliseconds including three launches and the transfers: 140 times
faster, the ratio of the bandwidths, because the job is \omterm{def:pl:accelerators-for-quantitative-work:bound}{memory-bound}. If 95\% of the overnight process is this kind of work and the rest
(reading positions, writing reports) stays on the host, Amdahl's law caps the whole process's gain at $1/(0.05 + 0.95/140) = 17.6$ times.
Against a server socket of many cores the per-core factor shrinks: \omterm{def:pl:accelerators-for-quantitative-work:bound}{memory-bound} work gains the ratio of the device's bandwidth to the
socket's, not to one core's.

\section{Training, and what the machine-learning book already said}

Book~12 (chapter~23) covered the accelerator's home ground: training neural networks is dominated by matrix products, \omterm{def:pl:accelerators-for-quantitative-work:bound}{compute-bound} with
high intensity, run in mixed precision on tensor units, and split across devices by data parallelism. The same roofline explains it: a
matrix product of size $n$ has intensity $n/12$, far above any ridge, and a device's reduced-precision peak is many times its
double-precision one. Quantitative finance's Monte Carlo is usually the other kind of kernel --- low intensity, double precision --- and
gains the bandwidth ratio, not the peak ratio.

\section{Deciding with a cost model}

A single intraday price request is small: ten thousand operations (a hundred cash flows on an interpolated curve) and a kilobyte crossing
the link (the trade and its result). On the core it takes 0.33~microseconds; on the device, 10.0 --- almost all of it the launch and
synchronisation latency, which the chapter takes as 10~microseconds (an assumption, not a datasheet figure). The device is 30 times slower.
Batched, the fixed cost is shared (\cref{fig:pl:accelerators-for-quantitative-work:batch}): the device breaks even at 31 trades per request
and approaches, for very large batches, 32.7 times the core --- a ceiling set by the link, because every trade's kilobyte still crosses it.

\begin{omfigure}
\begin{tikzpicture}
\begin{loglogaxis}[width=11.5cm,height=5.8cm,xlabel={trades per request (log scale)},ylabel={device speed-up (log scale)},
  xmin=1,xmax=1000000,ymin=0.02,ymax=100,grid=major,grid style={gray!20},tick label style={font=\small},label style={font=\small},
  table/col sep=comma]
\addplot[omDef,very thick] table[x=n,y=speedup]{figdata/platforms/14-accelerators-for-quantitative-work/batch.csv};
\draw[black,dashed] (axis cs:1,1) -- (axis cs:1000000,1);
\draw[omThm,thick] (axis cs:31,0.02) -- (axis cs:31,100) node[pos=0.12,right,font=\scriptsize]{break-even: 31 trades};
\node[font=\scriptsize,anchor=north east] at (axis cs:900000,28) {link-bound ceiling 32.7};
\end{loglogaxis}
\end{tikzpicture}

\omcaption{Predicted end-to-end speed-up of the device over one laptop core for a request pricing $n$ trades, with a launch and
synchronisation latency of 10~microseconds and the product page's PCIe bandwidth. A single trade is 30 times slower on the device; the device
wins from 31 trades, and no batch takes it beyond the link's ceiling. Data: \texttt{fig\_roofline.py}.}
\label{fig:pl:accelerators-for-quantitative-work:batch}
\end{omfigure}

\begin{example}[Faster overnight, not intraday]\label{ex:pl:accelerators-for-quantitative-work:verdict}
The same device, the same code style, two verdicts. The exposure job is large and \omterm{def:pl:accelerators-for-quantitative-work:bound}{memory-bound}, moves little across the link, and runs 140
times faster than on one core (17.6 times for the whole process by Amdahl's law). The price request is tiny, and the latency of a launch is
thirty times its whole cost on the CPU. The break-even batch depends on the latency assumed: 4 trades at 1~microsecond, 16 at 5, 31 at 10,
155 at 50. A pricing service that can batch its requests to hundreds per call would gain; one that answers each request as it comes cannot.
\end{example}

\begin{remark}[What the model leaves out]\label{rem:pl:accelerators-for-quantitative-work:limits}
The model is a bound: real kernels reach a fraction of either roof, as the table shows for the core. It ignores \omterm{def:pl:accelerators-for-quantitative-work:divergence}{thread divergence}, the cost of
porting and maintaining a second implementation, the device's price and power, and the queueing of requests on a shared device. It is enough
to rule an accelerator out, or to say that measuring one is worth the effort.
\end{remark}

\begin{dated}{2026-09}{A benchmark for pricing and risk hardware}\label{dat:pl:accelerators-for-quantitative-work:stac}
STAC Research describes its STAC-A2 benchmark suite as ``the industry standard for testing technology stacks used for compute-intensive
analytic workloads involved in pricing and risk management''. Vendors publish results on it for CPUs and accelerators; the chapter's roofline
is the back-of-the-envelope version of the same question.
\end{dated}

\section{Tutorial: two rooflines and a decision}\label{tut:pl:accelerators-for-quantitative-work:tut}

\begin{tutorial}
\textbf{Goal.} Measure one core's roof and four kernels on it, place them next to a cited device's roof, and predict two workloads end to
end. \textbf{End state:} \cref{fig:pl:accelerators-for-quantitative-work:roofline,fig:pl:accelerators-for-quantitative-work:batch}.
\begin{tutsteps}
\item \textbf{Measure}: \texttt{bench\_roofline.py} with one BLAS thread: the peak, the triad bandwidth, the four kernels.
\item \textbf{Place}: \texttt{Kernel.intensity} and \texttt{Roof.attainable} for each kernel on both roofs.
\item \textbf{Predict}: \texttt{overnight(cpu\_roof())}, \texttt{request(cpu\_roof(), n)}, \texttt{break\_even}.
\item \textbf{Stress the assumption}: \texttt{launch\_sensitivity}.
\end{tutsteps}
\textbf{What to change next.} Put the device's single-precision figures in the roof and see which workloads it changes; keep the market data
on the device between requests and send only the trade.
\end{tutorial}

\section{Build: the roofline model}\label{bld:pl:accelerators-for-quantitative-work:roofline}

\begin{build}
\bfield{Purpose} A cost model that tells the platform, before any purchase or port, which workloads an accelerator would speed up and by how
much, from measured and cited numbers.
\bfield{Interface} \texttt{Kernel(name, flops, bytes, transfer)}; \texttt{Roof(name, peak, bandwidth, link, launch, source)} with
\texttt{attainable}, \texttt{time}, \texttt{end\_to\_end}, \texttt{ridge}; \texttt{DEVICES}; \texttt{measure\_peak},
\texttt{measure\_bandwidth}; the reference kernels and their descriptors; \texttt{amdahl}; \texttt{break\_even}.
\bfield{Rules} Device figures are data with a source; CPU figures are measured on one core with one thread; a stated counting convention for
operations; assumptions (the launch latency) marked and swept.
\bfield{Acceptance tests} \texttt{code/firm/roofline/tests/}: the roof arithmetic; each reference kernel against a direct computation and its
counted intensity; Amdahl's law; the break-even search, including a device that never breaks even; the device data and the triad.
\bfield{Stretch} A measured roof of the whole socket; a device roof from a benchmark rather than a datasheet; the cost per result in money
from rental prices.
\end{build}

\begin{omsources}
\item S.~Williams, A.~Waterman and D.~Patterson, ``Roofline: an insightful visual performance model for multicore architectures'',
\emph{Communications of the ACM} 52(4), 2009.
\item G.~M. Amdahl, ``Validity of the single processor approach to achieving large scale computing capabilities'', AFIPS Spring Joint
Computer Conference, 1967.
\item NVIDIA, H100 product page; STAC Research, STAC-A2 benchmark page.
\end{omsources}

\section{Exercises}

\begin{exercise}[$\star$]\label{exo:pl:accelerators-for-quantitative-work:1}
What is the intensity of a matrix product of size $n$ under the chapter's convention, and above which $n$ is it \omterm{def:pl:accelerators-for-quantitative-work:bound}{compute-bound} on the core and
on the device?
\end{exercise}

\begin{exercise}[$\star$]\label{exo:pl:accelerators-for-quantitative-work:2}
Why is the exposure job's speed-up close to the bandwidth ratio and not to the peak ratio?
\end{exercise}

\begin{exercise}[$\star$]\label{exo:pl:accelerators-for-quantitative-work:3}
With 95\% of a process sped up 140 times, what is the whole process's gain, and what would it be with 99\%?
\end{exercise}

\begin{exercise}[$\star\star$]\label{exo:pl:accelerators-for-quantitative-work:4}
Why does the path step reach only 16\% of its roof, and what would move it closer?
\end{exercise}

\begin{exercise}[$\star\star$]\label{exo:pl:accelerators-for-quantitative-work:5}
Where does the ceiling of 32.7 for large batches of price requests come from?
\end{exercise}

\begin{exercise}[$\star\star$]\label{exo:pl:accelerators-for-quantitative-work:6}
Which parts of an American-option Monte Carlo with a regression at each date would suffer from \omterm{def:pl:accelerators-for-quantitative-work:divergence}{thread divergence}, and which would not?
\end{exercise}

\begin{exercise}[$\star\star\star$]\label{exo:pl:accelerators-for-quantitative-work:7}
\emph{Coding.} With \texttt{firm.roofline}, recompute the break-even batch if the market data (1~MB) must be sent with every request instead
of staying on the device.
\end{exercise}

\begin{exercise}[$\star\star\star$]\label{exo:pl:accelerators-for-quantitative-work:8}
\emph{Find the flaw.} ``The device has 508 times our core's peak, so our risk run, which takes ten hours on one core, will take a minute.''
\end{exercise}

\section{Problem: Faster Overnight, Not Intraday}

\begin{problem}[{Weekend problem --- two workloads, one device}]\label{pb:pl:accelerators-for-quantitative-work:1}
One laptop core as measured, the device of \cref{dat:pl:accelerators-for-quantitative-work:h100}, the chapter's two workloads.

\medskip\noindent\textbf{Part I --- Roofs.}
\begin{enumerate}
\item What are the core's peak, bandwidth and ridge?
\item What are the device's, and the ratios to the core?
\item Which of the four kernels are \omterm{def:pl:accelerators-for-quantitative-work:bound}{memory-bound} on the core, and on the device?
\item Why does the path step sit below its roof?
\item What does a roofline not capture?
\end{enumerate}

\medskip\noindent\textbf{Part II --- The overnight job.}
\begin{enumerate}[resume]
\item How many operations and bytes does the exposure job have, and what crosses the link?
\item What are its times on the core and on the device?
\item Why is the speed-up what it is?
\item What does Amdahl's law make of it for the whole process?
\item How would the comparison change against a many-core server?
\end{enumerate}

\medskip\noindent\textbf{Part III --- The request.}
\begin{enumerate}[resume]
\item What are a single request's times on the core and on the device?
\item At what batch does the device break even?
\item How does the break-even move with the launch latency?
\item What bounds the speed-up of very large batches?
\item What would a pricing service have to change to benefit?
\end{enumerate}

\medskip\noindent\textbf{Part IV --- The verdict.}
\begin{enumerate}[resume]
\item State the \emph{named result}: the predicted end-to-end speed-up of the exposure Monte Carlo and of a single-trade price request, and
the batch size at which the device breaks even.
\item Which workloads of the platform would you move to a device first?
\item What would you measure before buying one?
\item Why use double precision in the model?
\item In one sentence: when does an accelerator pay?
\end{enumerate}
\end{problem}

\section{Interview questions}

\begin{interviewq}[$\star$ \iqroles{developer}]\label{iq:pl:accelerators-for-quantitative-work:1}
What is \omterm{def:pl:accelerators-for-quantitative-work:intensity}{arithmetic intensity}, and why does it matter for choosing hardware?
\end{interviewq}

\begin{interviewq}[$\star\star$ \iqroles{developer, researcher}]\label{iq:pl:accelerators-for-quantitative-work:2}
Your Monte Carlo got 30 times faster on a graphics processor, your pricing service got slower. Explain.
\end{interviewq}

\begin{interviewq}[$\star\star$ \iqroles{developer}]\label{iq:pl:accelerators-for-quantitative-work:3}
How would you estimate, without the hardware, what a device would do for a workload?
\end{interviewq}

\begin{interviewq}[$\star\star$ \iqroles{researcher}]\label{iq:pl:accelerators-for-quantitative-work:4}
Can you run a Monte Carlo on a device and get exactly the CPU's numbers? How?
\end{interviewq}

\begin{interviewq}[$\star\star\star$ \iqroles{developer}]\label{iq:pl:accelerators-for-quantitative-work:5}
Design the decision process for moving parts of a risk platform to accelerators.
\end{interviewq}
